\chapter{催化剂怎样另辟蹊径}

\begin{kaichang}
你家药箱里可能有一小瓶棕色的双氧水——擦破皮时，倒一点在伤口上，立刻冒起一层白沫，嘶嘶作响。理发店里染头发用的也是它，只是更浓一些。

现在做两个思想实验。第一，把这瓶双氧水放在柜子里，半年后再看：它还是清清亮亮的一瓶水，几乎没有变化。第二，切一小块生猪肝，蘸一点同样的双氧水：剧烈的气泡翻滚而出，像烧开了一样。

同一瓶液体，为什么自己待着安安静静，一碰到猪肝就"发疯"？猪肝给了双氧水什么——能量？指令？还是别的什么东西？猪肝自己又付出了什么代价，它被"用完"了吗？把猜测写下来，这一章结束时回来对答案。答案藏在一个贯穿工业、生命和你身体每个细胞的词里：\textbf{催化剂}。
\end{kaichang}

\section{一瓶"慢性子"的消毒水}

双氧水是过氧化氢（\chem{H_2O_2}）的水溶液，药店里卖的大约只有 3\% 的浓度。过氧化氢分子比水分子只多一个氧原子，可就是这一个多出来的氧，让整个分子坐立不安：它随时可能一分为二，变成水和氧气。这个变化放出热量，从能量上看完全是"下坡"——用第 28 章的话说，它的 $\Delta G$ 是负的，属于热力学允许甚至鼓励的自发变化。

但"自发"不等于"快"。第 28 章反复强调过这一点，双氧水是最好的活教材：一瓶纯净的双氧水在室温下放上一年，分解掉的也只有一点点。它在能量上"想"往下走，却被一堵高墙挡着——第 29 章说过，这堵墙叫\textbf{活化能}。过氧化氢分子要分解，先得把分子扭成一个高能量的别扭形状（过渡态），常温下几乎没有什么碰撞能提供这么大的能量，于是绝大多数分子只能原地待着。

现在把那小块猪肝放进去。几秒钟之内，杯子里泡沫翻腾，杯壁微微发热，伸一根带火星的木条进去（这步交给大人），木条会重新燃起来——那是氧气。同样的现象你早就见过：双氧水滴在伤口上冒泡，正是因为血液和组织里含有催化这个分解的蛋白质。

再看几个生活里的同类事件：

\begin{itemize}[itemsep=2pt]
\item 白米饭或馒头嚼久了会发甜，因为唾液里的淀粉酶在把淀粉拆成糖。米饭在嘴里放着不嚼，甜味来得慢得多。
\item 汽车排气管中段有一个不起眼的金属罐子，里面涂着铂、铑、钯等贵金属，尾气里的一氧化碳和氮氧化物经过它就被转化成二氧化碳和氮气。没有它，每辆车都是一座移动毒气工厂。
\item 有一种白金（铂）催化打火机，不用火石：按一下，燃料气体喷到铂网上就自行发热起火。铂网摸上去完好无损，却天天在"点火"。
\item 氮肥厂里，氮气和氢气在一层层铁颗粒上流过，变成氨。没有这道工序，全世界的农田养不活今天一半的人口。
\end{itemize}

还有一类催化剂你可能天天穿在身上：加酶洗衣粉里的蛋白酶和脂肪酶，专门拆掉奶渍、血渍和油渍——它们只能用温水，开水一烫就失效，原因这一章末尾你就明白了。冬天有人揣着"白金怀炉"取暖：煤油蒸气碰到铂催化剂，不用明火就持续氧化放热，摸上去温温的，能热一整天。从洗衣粉到怀炉，催化剂早就混进了日常，只是它干活不留名。

这些场合有一个共同点：某种物质让反应变快了，而它自己似乎没有被消耗。这样的物质叫\textbf{催化剂}。先别急着背定义，我们用一个小实验把"没有消耗"看清楚。

\begin{shiyan}
\textbf{材料}：药店买的 3\% 双氧水一小瓶；一小块生猪肝（或生土豆、一小撮酵母粉，任选其一）；两个透明杯子；厨房秤（可选）；筷子。

\textbf{步骤}：两个杯子各倒约三分之一杯双氧水。一杯留作对照，另一杯投入一小块猪肝（约指甲盖大），用筷子压住观察。十分钟后，把猪肝捞出来冲一冲，和原来比一比大小、颜色。如果你用了厨房秤，可以反应前后各称一次猪肝块（吸干水分再称）。

\textbf{观察点}：①两杯冒泡速度的差别——对照杯几乎不冒泡，猪肝杯泡沫翻滚；②摸猪肝杯的外壁，是不是温热的？放热说明反应确实发生了，能量账目没变；③捞出的猪肝变小了吗？它并没有被"吃掉"。

\textbf{安全提示}：只用 3\% 的医用双氧水，浓度更高的（如理发店用或工业用）会灼伤皮肤，绝对不许碰。双氧水别溅进眼睛，操作时最好戴手套。带火星木条复燃的验证涉及明火，只能看老师演示或视频。实验后的猪肝不能吃了，直接丢弃。
\end{shiyan}

猪肝没有被消耗，它只是"路过"了一下，双氧水就从慢性子变成了急脾气。它是怎么做到的？

\section{不是推一把，而是换条路}

把镜头推进到分子层面。第 29 章给过一幅图景：反应物分子要变成产物，必须先翻越一座能量山头，山顶上那个高能量的别扭状态叫过渡态。山越高，能翻过去的分子越少，反应就越慢。过氧化氢自己分解，要翻的就是一座相当高的山：分子里的化学键得先被拉到快要断开的程度，单凭室温下的随机碰撞，百万次、亿次里也难得成功一次。

猪肝里起作用的，是一种叫\textbf{过氧化氢酶}的蛋白质；化学老师演示用的则是二氧化锰（\chem{MnO_2}）的黑色粉末。两者干的活本质相同：\textbf{它们不是在山后面推分子一把，而是给分子指了另一条路——一条分成几段、每段山都不高的路。}

以二氧化锰为例，简化地说，过程分两步。第一步，一个过氧化氢分子撞到二氧化锰表面，被表面的锰原子"接住"，与它形成暂时的结合，原来的旧键在这个过程中被削弱、拆开，生成水和一些附着在表面上的高活性碎片。第二步，另一个过氧化氢分子过来，和这些碎片结合，氧原子两两配对成氧气分子离开，二氧化锰表面恢复原样，等着接下一批客人。

请注意这里最要紧的一点：催化剂\textbf{并不是置身事外的旁观者}。它真实地和反应物牵了手，形成了只存在一瞬间的"中间体"，然后在下一段路上体面地退出，毫发无损地回到起点。化学家把这叫作一个"催化循环"：进去，干活，出来，再进去。一粒二氧化锰、一个过氧化氢酶分子，一生中可以重复这个循环成千上万次乃至上亿次——这就是一小撮粉末、一小块猪肝能"搞定"一整杯双氧水的原因。

下面这幅能量图把"换路"画了出来。注意：这是人为的表示法，曲线不代表任何真实形状，只表示能量高低。

\begin{center}
\begin{tikzpicture}[scale=0.85]
\draw[->] (0,0) -- (10,0) node[below] {反应进程};
\draw[->] (0,0) -- (0,4.6) node[above] {能量};
\draw[thick] (0.4,2.6) .. controls (2.2,2.6) and (2.4,4.3) .. (3.4,4.3)
  .. controls (4.4,4.3) and (4.6,1.0) .. (6.4,1.0);
\draw[thick,dashed] (0.4,2.6) .. controls (1.4,2.6) and (1.5,3.1) .. (2.2,3.1)
  .. controls (2.9,3.1) and (3.0,1.9) .. (3.8,1.9)
  .. controls (4.6,1.9) and (4.7,3.1) .. (5.3,3.1)
  .. controls (5.9,3.1) and (5.6,1.0) .. (6.4,1.0);
\draw[<->] (7.6,1.0) -- (7.6,4.3) node[midway,right] {原路的山};
\draw[<->] (8.6,1.0) -- (8.6,3.1) node[midway,right] {新路的山};
\node at (1.0,2.85) {反应物};
\node at (6.4,0.6) {产物};
\node[align=center] at (3.4,5.0) {实线：无催化\quad 虚线：有催化};
\end{tikzpicture}
\end{center}

虚线把一座高山换成了几座矮山。矮山里最矮的那座，也比没有路强——常温下就能有可观的分子翻过去，反应于是成千上万倍地加快。而起点（反应物）和终点（产物）的高度，两条路完全一样。这一点，下一节要掰开揉碎地讲，因为它是全世界化学考试和化学骗局共同的"重灾区"。

\section{快，和"多"，是两回事}

现在正面回答一个最容易被搞混的问题：催化剂到底改变了什么，没改变什么？

先看没改变的。第 27 章讲过赫斯定律：一个反应的焓变 $\Delta H$ 只看起点和终点，与路径无关。双氧水分解成水和氧气放出多少热，无论你让它慢慢自己分解、用二氧化锰催，还是用猪肝催，放出的总热量一模一样——起点和终点没变，账本就一个字都不会改。同理，第 28 章的 $\Delta G$ 也只看始末状态：方向由它定，极限由它定，催化剂碰不到它。

还有平衡。第 31 章会专门讲化学平衡，这里先记住结论：一个可逆反应到达平衡时停在什么位置（反应物和产物各占多少），由平衡常数 $K$ 决定，而 $K$ 的数值完全由能量账目决定。催化剂\textbf{不改变平衡常数，不改变平衡位置}。它做的事只有一件：让体系\textbf{更快到达}那个本来就注定的位置。原因在于催化剂对正反应和逆反应一视同仁——它给正向修的路，反方向同样好走。正向快多少倍，逆向也快多少倍，两者的比值不变，终点自然不变。

打个比方：山脚是反应物，山顶另一边的低谷是产物。两地的高度差是老天爷定的。原来只有一条翻越主峰的险路，人迹罕至；催化剂就像新修的穿山隧道，把通行时间从几天缩到几小时。可是隧道改变不了山脚和低谷的高度差，也改变不了"住在低谷更舒服"这个事实——它只是让两边的人往来更快而已。发洪水时隧道两头都会被淹，方向还是取决于哪里地势低。

\begin{qiaoliang}
催化剂能把反应加快多少？第 29 章给过阿伦尼乌斯关系的直觉：速率对活化能 $E_a$ 极其敏感，大致按 $e^{-E_a/(RT)}$ 缩放，其中 $R=\ud{8.314}{J\cdot mol^{-1}\cdot K^{-1}}$，$T$ 是热力学温度。假设某反应在室温 $298\ \mathrm{K}$ 下，活化能从 $\ud{75}{kJ\cdot mol^{-1}}$ 被催化剂降到 $\ud{50}{kJ\cdot mol^{-1}}$，速率提升的倍数是
\[
\frac{k_{\text{催}}}{k_{\text{原}}} = e^{(75-50)\times 10^{3}/(8.314\times 298)} = e^{25000/2478} \approx e^{10.1} \approx 2.4\times 10^{4}.
\]
区区 $\ud{25}{kJ\cdot mol^{-1}}$ 的"削山"——大约相当于一根普通共价键能量的零头——就把反应加快了两万多倍。这就是为什么指数函数是化学里最有威力的数学：活化能只差一点点，速度就差出几个数量级。肝里的过氧化氢酶更夸张，它把这个分解反应的速率提高了约 $10^{14}$ 倍，一个酶分子每秒钟最多能处理上百万个过氧化氢分子。你肝脏每天代谢产生的过氧化氢，如果不及时清除会损伤细胞，而一小把酶分子就把它们秒速拆成了水和氧气。
\end{qiaoliang}

\section{符号怎么记这笔账}

到现在，才轮到符号登场。催化剂在化学方程式里的写法很特别：它不进反应物，也不进产物，而是写在\textbf{箭头的上方}，像一个"路况标注"：

\[\chem{2H_2O_2} \xrightarrow{\chem{MnO_2}} \chem{2H_2O} + \chem{O_2}\]

这个写法本身就是化学家的诚实：二氧化锰没有出现在账本的任何一栏里，它只是这条路的使用条件。你身体的细胞里，同一个反应由过氧化氢酶完成，写法一样：

\[\chem{2H_2O_2} \xrightarrow{\text{过氧化氢酶}} \chem{2H_2O} + \chem{O_2}\]

再看两个将要反复见面的反应。哈伯法合成氨（第 31 章还会回来细算它的平衡账）：

\[\chem{N_2} + \chem{3H_2} \rightleftharpoons \chem{2NH_3}\]

铁做催化剂，写在可逆箭头上方；注意这是个可逆反应，铁同时加速了向左和向右两个方向。汽车尾气催化器里发生的事之一：

\[\chem{2CO} + \chem{2NO} \rightarrow \chem{2CO_2} + \chem{N_2}\]

铂、铑等贵金属是这条路的"隧道施工队"。顺便说一句，这个反应的 $\Delta G$ 本来就是负的——一氧化碳和一氧化氮共处时，变成二氧化碳和氮气在能量上完全划算。它们之所以能在尾气里"和平共处"一路排出，纯粹是因为没有路可走；催化器做的全部工作，就是把那条路修通。

\section{三条不同的"修路方案"}

催化剂修路的手段并不止一种，常见的可以分成三类。

\textbf{固体表面催化}（多相催化）。二氧化锰、铂网、合成氨的铁颗粒都是这一类。反应物分子吸附到固体表面，被表面的原子拉住、固定、弯折，旧键在"按压"下变弱，新键在邻居的撮合下生成，产物随后脱附离开。因为活只在表面上干，表面积就是命根子：工业催化剂都被做成多孔的颗粒或蜂窝，一克铂催化剂展开的内表面积可以有一个房间那么大。这也解释了为什么催化剂怕灰尘、怕油污——表面一旦被占，路就废了。

\textbf{酸碱催化}（均相催化的一种）。氢离子 \chem{H^+} 是个身手敏捷的"中转站"：它先把自己塞给反应物某个原子，让某根键变得脆弱易断；等反应走完，氢离子又被原样吐出来。糖水解成葡萄糖和果糖、淀粉在稀酸里煮成糊精，都是氢离子在循环打工。酯化反应反过来用酸催化也是同一个道理（第二册有机化学里还会再见）。

\textbf{酶催化}。酶是蛋白质折叠成的精密口袋，口袋的形状、电荷、亲疏水性都经过上亿年进化的打磨，恰好抱住某一种（或一类）分子，把它摆成最容易起变化的姿势。过氧化氢酶、淀粉酶、你消化道的每一种消化酶都是这一类。酶的效率高、选择性强，但也有蛋白质的娇气：温度高了、酸碱不对了，口袋一变形就罢工。酶的世界太精彩，也太重要，第二册第 49 章会用整整一章来讲它，这里只记住一点：酶\textbf{同样是催化剂}，同样遵守这一章的全部规则——不提供能量，不改平衡，只换条路。

\section{哈伯的赌局：给化学反应修路的人}

\begin{zhengju}
二十世纪初，全世界面临一场静悄悄的危机：农田缺氮。氮是蛋白质和遗传物质的骨架元素，空气里明明五分之四都是氮气，可氮分子里的三根键牢得出奇（第 17 章讲过这根键有多结实），植物根本用不上。当时农业依赖智利硝石和鸟粪，储量肉眼可见地在枯竭。科学家算过一笔账：如果不解决氮肥问题，按当时的人口增长，几十年内就会出现大规模饥荒。

氮气和氢气合成氨在热力学上是可行的——问题全在速度。氮分子那根三键意味着高得吓人的活化能，没有催化剂时，这个反应慢到等于不存在。1909 年，德国化学家哈伯用锇和铀做催化剂，在高温高压下第一次让氨以看得见的速度流出。但真正把这个反应变成工业的，是巴斯夫公司的博施和他的团队。锇太贵太稀有，博施的助手米塔希做了一件笨到伟大、也伟大到辉煌的事：系统地试验了上万种催化剂配方——据记载超过两万个实验、三千多种候选材料——最后发现最便宜的铁加上一点助剂，效果反而最好。

这条路一修通，人类从此可以"从空气里烤出面包"。今天，全世界生产的氨养活的粮食大约支撑了地球一半人口的蛋白质摄入；这套装置同时消耗着全球约 1\% 到 2\% 的能源——为降低这点能耗，化工工程师至今仍在和催化剂较劲。而硬币的另一面同样真实：氨也是炸药的原料，哈伯后来又主导了毒气战的研究。催化剂本身没有立场，它只是把路修通；走向肥料还是走向炸药，是握着方向盘的人的选择。这个判断，化学给不了，得你自己给。
\end{zhengju}

\section{催化剂也会"生病"}

催化剂不是永动机，它有自己的脆弱，这些脆弱反过来塑造了半个工业世界。

\textbf{中毒}。如果某种杂质分子特别爱粘在催化剂的活性位点上，赖着不走，路就被堵死了，这叫催化剂中毒。最有名的连锁反应是：硫化物会让合成氨的铁催化剂中毒，所以原料气必须深度脱硫；铅会让汽车催化器的铂中毒，所以上世纪七八十年代起，各国陆续把汽油里的含铅添加剂赶了出去——一场影响深远的"无铅化"，最初的推动力之一竟然是为了保护排气管里那几克贵金属。燃料电池里的铂也怕一氧化碳，微量的一氧化碳就能让电极"窒息"，这是氢燃料必须高纯的原因之一。

\textbf{选择性}。同一批反应物，往往有好几条能量上都说得通的路。好的催化剂不只快，还要"专一"——只修通向目标产物的那条。合成气的碳一化学、制药厂里不对称合成手性分子（第 44 章会讲手性为什么关乎药效），靠的都是催化剂的选择性。酶把选择性做到了极致：过氧化氢酶对过氧化氢一心一意，几乎不碰别的分子。

\textbf{循环与再生}。因为催化剂反应后复原，它原则上可以无限循环——这正是它"省"的本质。但现实中的催化剂会慢慢烧结、积碳、中毒，于是工厂发展出各种"再生"手艺：烧焦积碳、化学清洗、重新分散。炼油厂的流化催化裂化装置里，催化剂颗粒在反应器和再生器之间日夜不停地循环，像一条流动的传送带，烧掉结焦，满血复活，再回去干活。

\section{这副地图在哪里失效}

任何模型都有边界，催化剂的边界同样清晰，而且每一条边界上都挂着真实发生过的骗局。

第一，催化剂\textbf{救不活热力学上死掉的反应}。$\Delta G$ 为正在给定条件下走不通的方向，修再快的路也是白修——路上车水马龙，但两头都不挪窝。上世纪轰动一时的"水变油"闹剧，违背的就是这条：水里没有碳，能量账也对不上，无论加什么"神奇催化剂"都不可能变成汽油。今后凡听到"某某催化剂让不可能的反应发生了"，你可以直接判定：要么是测量错了，要么是有人在撒谎。

第二，催化剂\textbf{不改产量上限}。它只改变到达平衡的快慢。合成氨工厂里，想把氨的产率从三成提到五成，靠的不是换更强的催化剂，而是改温度、改压强、及时把氨移走——那些是平衡的语言，第 31 章细讲。

第三，催化剂\textbf{自己有工作条件}。温度太低，吸附不上；温度太高，酶变性、金属颗粒烧结长大、表面积塌掉。每个催化剂都有自己舒服的"工作窗口"，这正是工业反应器设计最费心思的部分。

\begin{wuqu}
"催化剂不参与反应，所以反应前后它没有任何变化。"——这句话半对半错，错的半句流传更广。催化剂当然\textbf{参与}反应：它吸附反应物、形成中间体、再把自己释放出来，只是最后总量不变。而且连"没有任何变化"也常常不成立：用过的铂网表面会变得粗糙，汽车催化器里的贵金属会慢慢迁移聚集，酶分子有一定的寿命会老化失活。准确的说法是：催化剂在\textbf{化学计量上}不变——它不进总账；但在\textbf{物理状态和个体寿命上}，它和你我一样会被岁月磨损。它换来的是速度，付出的是慢慢老去。
\end{wuqu}

\section{回到那块猪肝}

现在可以对答案了。猪肝里有过氧化氢酶，它的活性口袋中心坐着一个铁原子（正是第 10 章说的、过渡金属的拿手好戏）。双氧水分子钻进口袋，铁原子先"接"走一个氧，再把它交给下一个双氧水分子，自己恢复原状——一个催化循环，每秒重复上百万次。气泡是氧气，白沫是被气泡鼓起来的液体，微热是分解放出的热。

猪肝被"用完"了吗？没有。猪肝变小了，是因为组织被搅动、被破坏，而不是酶被消耗。酶分子反应后原样还在，随时可以接下一个活——事实上，一块煮过的猪肝就不灵了，因为高温把蛋白质的折叠口袋烫散了架（这个"变形就失活"的现象，叫蛋白质的变性，第二册会专门讲）。猪肝给双氧水的不是能量，也不是指令，而是一条\textbf{现成的、低矮的、精确铺好的路}。

顺带说，生土豆、酵母也能干这件事，只是它们的过氧化氢酶含量不同，冒泡的凶猛程度不同——你实验里如果换了材料，观察到的差别就是催化能力的差别。

\section{这条路通向生命}

这一章你认识了化学世界最了不起的一类角色：不改变规则，却在规则之内把速度玩弄于股掌之间的"修路者"。没有催化剂，许多热力学上完全允许的反应慢得等于不存在；有了催化剂，"能发生"和"来得及发生"之间那道鸿沟被填平了。

而生命，恰恰是这场修路艺术的最大受益者。你的每一个细胞里都同时进行着上千个反应，它们中的绝大多数如果离开酶，在体温下慢得可以忽略不计——换句话说，\textbf{生命就是一个被催化精确调度着的化学反应网络}。酶凭什么比工业催化剂更高效、更专一？细胞又怎样通过控制酶的数量和活性，来决定哪个反应开、哪个反应关？这条线将把我们一路带进第二册：第 49 章专门解剖酶这台分子机器，而后的代谢、呼吸、光合作用，每一步都是催化剂的舞台。从下一章开始，我们先补齐最后一块拼图——当反应快起来之后，它会在哪里停下来：化学平衡。

\begin{tiaozhan}
为什么催化剂加速正反应，就必然以同样的倍数加速逆反应？这背后是一条深刻的对称性，叫微观可逆性：正反应翻越的过渡态，和逆反应翻越的是\textbf{同一个山顶}。设原路正、逆反应的活化能分别是 $E_a^{+}$ 和 $E_a^{-}$，产物比反应物低 $\Delta H$，则 $E_a^{+}-E_a^{-}=\Delta H$。催化剂把山顶削低同一个数值 $\Delta E$，正、逆活化能就同时减小 $\Delta E$，差值不变——于是正、逆速率被放大同样的倍数 $e^{\Delta E/(RT)}$，比值 $k^{+}/k^{-}$ 纹丝不动。而平衡常数 $K$ 正比于这个比值，所以 $K$ 不变。结论干净利落：\textbf{凡能加速到达平衡的物质，必然不改平衡本身}。这不是经验总结，是能量几何的必然结果。如果有人声称发明了"只加速正反应"的催化剂，那和永动机是同一类东西。
\end{tiaozhan}

\section*{练一练}

\begin{sikaoti}
\item 画一幅能量图：横轴反应进程、纵轴能量，画出同一反应"无催化"和"有催化"两条路径。用箭头标出：两种情况的活化能各是哪一段？$\Delta H$ 是哪一段？它在两条路径里相同吗？
\item 合成氨工厂里，工程师把铁催化剂换成效率翻倍的新催化剂。请问：氨的\textbf{平衡产率}会变吗？工厂的\textbf{日产量}可能怎么变？用"快"和"多"两栏分别作答。
\item 有人宣称发明了一种催化剂，能让水在常温常压下分解成氢气和氧气，不需要任何能量输入。用水电解的能量账说明这为什么不可能。他提供的"氢气"更可能来自哪里？
\item 设计一个实验，检验二氧化锰在双氧水分解前后是否真的"没有被消耗"。提示：不溶的黑色固体可以过滤、洗涤、烘干、称量。还要排除什么干扰？
\item 含铅汽油为什么被全球淘汰？从"催化剂中毒"的角度写出这条因果链：含铅添加剂 $\rightarrow$ \dots $\rightarrow$ 尾气污染加重。
\item 猪肝煮过之后催化能力大幅下降，而二氧化锰烧红后依然有效。两类催化剂的这个差别，根源在它们各自的"修路工具"是什么？
\end{sikaoti}
